\section{实证结果报告模板}

本节给出正式实证结果的写作模板。当前仓库尚未导入正式样本，因此以下表格以占位形式展示。待运行 \texttt{src/pipeline/run\_regression.py} 生成正式结果后，可直接替换相应数值与叙述。

\subsection{描述性统计与相关性分析}

正式写作时，可先报告样本企业在真实性指数、包装性指数和经营结果变量上的总体分布，说明是否存在明显异质性。一个可直接替换的写法如下：

\begin{quote}
描述性统计结果显示，样本企业 AI 披露真实性指数存在明显差异，说明 IPO 企业在数字化和 AI 叙事上并非同质化表达。上市后经营表现指标也呈现较大离散度，为后续回归分析提供了基础。相关系数结果初步表明，真实性指数与 ROA、营业收入增长率和总资产周转率大体呈正相关，而包装性指数与上述经营结果变量总体呈负相关，但具体关系仍需在控制变量和固定效应条件下进一步检验。
\end{quote}

\begin{table}[H]
\centering
\caption{描述性统计表模板}
\begin{tabular}{lcccc}
\toprule
变量 & 均值 & 标准差 & 最小值 & 最大值 \\
\midrule
真实性指数 $AuthIndex$ & 待填 & 待填 & 待填 & 待填 \\
包装性指数 $PackagingIndex$ & 待填 & 待填 & 待填 & 待填 \\
ROA & 待填 & 待填 & 待填 & 待填 \\
营业收入增长率 & 待填 & 待填 & 待填 & 待填 \\
总资产周转率 & 待填 & 待填 & 待填 & 待填 \\
企业规模 & 待填 & 待填 & 待填 & 待填 \\
资产负债率 & 待填 & 待填 & 待填 & 待填 \\
研发强度 & 待填 & 待填 & 待填 & 待填 \\
\bottomrule
\end{tabular}
\end{table}

\subsection{基准回归结果}

基准回归部分建议先分别报告真实性指数对三个经营结果变量的系数，再同时纳入包装性指数。若正式回归结果符合预期，可采用如下文字模板：

\begin{quote}
基准回归结果表明，AI 披露真实性指数对企业上市后经营表现具有显著正向影响。以 ROA 为例，真实性指数的回归系数为 [待填]，在 [待填] 水平上显著；以营业收入增长率和总资产周转率为被解释变量时，真实性指数的系数同样为正，说明 IPO 阶段更加具体、审慎且具有信息含量的 AI 披露，更可能在上市后兑现为盈利能力、成长性和运营效率改善。相较之下，包装性指数的系数整体为负，表明概念化叙事越强、兑现难度越大。
\end{quote}

\begin{table}[H]
\centering
\caption{基准回归结果模板}
\resizebox{\textwidth}{!}{
\begin{tabular}{lcccc}
\toprule
变量 & ROA & 营业收入增长率 & 总资产周转率 & 综合经营表现指数 \\
\midrule
真实性指数 $AuthIndex$ & 待填 & 待填 & 待填 & 待填 \\
 & (待填) & (待填) & (待填) & (待填) \\
包装性指数 $PackagingIndex$ & 待填 & 待填 & 待填 & 待填 \\
 & (待填) & (待填) & (待填) & (待填) \\
企业规模 & 已控制 & 已控制 & 已控制 & 已控制 \\
资产负债率 & 已控制 & 已控制 & 已控制 & 已控制 \\
企业年龄 & 已控制 & 已控制 & 已控制 & 已控制 \\
研发强度 & 已控制 & 已控制 & 已控制 & 已控制 \\
所有制性质 & 已控制 & 已控制 & 已控制 & 已控制 \\
行业固定效应 & 是 & 是 & 是 & 是 \\
年份固定效应 & 是 & 是 & 是 & 是 \\
样本量 & 待填 & 待填 & 待填 & 待填 \\
$R^2$ & 待填 & 待填 & 待填 & 待填 \\
\bottomrule
\end{tabular}
}

\vspace{0.3em}
\parbox{0.92\textwidth}{\small 注：括号内为稳健标准误。显著性水平可按 *、**、*** 标注。}
\end{table}

\subsection{维度拆分与稳健性检验}

在基准回归之后，建议进一步检验三个维度的作用差异。预计场景具体性和表述审慎性会比单纯提及强度更能稳定预测上市后经营表现，因为“说得多”不一定代表“做得实”，而“说得具体、说得克制”更可能反映真实能力基础。若结果支持这一判断，则可用如下表述：

\begin{quote}
分维度回归结果显示，提及强度本身并不稳定预测企业上市后经营表现，而场景具体性和表述审慎性具有更强解释力。这说明企业是否真正把 AI / 数字化叙事落实到业务流程、系统建设和组织实践中，才是区分真实转型与概念包装的关键。
\end{quote}

此外，本文还可以进行以下稳健性检验：一是将被解释变量替换为综合经营表现指数；二是将时间窗口扩展到上市后第二年；三是剔除极端值或高波动行业样本；四是将包装性指数单独作为核心解释变量。若回归结果在上述检验中保持方向一致，则可以进一步增强结论可信度。

\begin{table}[H]
\centering
\caption{稳健性检验结果模板}
\begin{tabular}{lccc}
\toprule
检验项目 & 核心变量 & 系数方向 & 结论 \\
\midrule
替换为综合经营表现指数 & 真实性指数 & 待填 & 待填 \\
上市后第二年窗口 & 真实性指数 & 待填 & 待填 \\
剔除极端值样本 & 真实性指数 & 待填 & 待填 \\
包装性单独回归 & 包装性指数 & 待填 & 待填 \\
\bottomrule
\end{tabular}
\end{table}

\subsection{结果讨论}

若正式结果显示真实性指数显著为正、包装性指数显著为负，则可以从以下逻辑展开讨论。第一，IPO 阶段的 AI 叙事并非完全是资本市场热点附着，它在一定程度上确实反映企业数字化能力和治理质量差异。第二，真正有信息含量的数字化披露并不主要表现为高频口号，而是表现为具体业务场景、系统建设证据和边界清晰的表述。第三，问询回复作为监管互动文本，为识别企业是否存在概念包装提供了重要补充，因此将其与招股说明书联合使用具有方法上的必要性。

如果部分结果不显著，也可以据此讨论 IPO 场景中数字化叙事的“高频但异质”特征，即企业提及 AI 并不自动转化为经营改善，真正起作用的可能只是其中更具体、更可验证的那一部分表述。
